%& -shell-escape
% !TeX root = thesis-abstract.tex
% !TeX encoding = UTF-8
% !TeX program = XeLaTeX

\documentclass{.local/lib/classes/ndvr-super-article-standalone}

\begin{document}

    Таким образом, сформулирован алгоритм индексации и поиска
    нечётких дубликатов видео:
    \begin{itemize}
        \item для испытуемого потока видео:
        \begin{enumerate}
            \item выделить аномалии видео (границы съёмок) 
                с~помощью сравнения скользящих 
                средних оценок сигнала;
            \item построить дескрипторы съёмок:
            \begin{enumerate}
                \item вычислить относительные длины съёмок,
                \item вычислить характеристики граничных кадров;
            \end{enumerate}
            \item найти дескрипторы съёмок в индексе
                  с помощью модифицированного алгоритма $\iSAX{2.0}$:
            \begin{enumerate}
                \item выполнить кусочно-агрегатную аппроксимацию
                      дескрипторов съёмок,
                \item выполнить поиск по индексному дереву;
            \end{enumerate}
            \item если дескрипторы съёмок искомого видео найдены:
            \begin{itemize}
                \item сравнить характеристики 
                      граничных кадров съёмок,
                \item в случае успеха 
                      пометить искомое видео как дубликат;
            \end{itemize}
            \item в противном случае: 
            \begin{enumerate}
                \item добавить испытуемое видео 
                    в индекс с помощью модифицированного 
                    алгоритма $\iSAX{2.0}$,
                \item перестроить индексное дерево.
            \end{enumerate}
        \end{enumerate}
    \end{itemize}

    %\pagebreak

    \noindent
    На рисунке \ref{fig:video-search-dfd} 
    звёздочками отмечены элементы, 
    участвующие при реализации 
    сформулированного алгоритма.


    \begin{figure}[p]
        \includegraphics[width=\textwidth]
            {.local/vec/out/pdf/video-model-dfd-state.pdf}
        \caption{DFD-диаграмма поиска по видео 
            в нотации Йордона — Де Марко.}
        \label{fig:video-search-dfd}
    \end{figure}

    В отличие от методов, рассмотренных ранее,
    алгоритм обладает достоинствами:
    \begin{itemize}
        \item алгоритм работает в режиме реального времени;
        \item алгоритм может быть применен 
        к видео различных типов без полного 
        перестроения индекса;
        \item благодаря тому, что при поиске не учитываются 
        пространственные характеристики видео, 
        алгоритм неприхотлив к вычислительным ресурсам.
    \end{itemize}

    Предложенный алгоритм обладает рядом недостатков:
    \begin{itemize}
        \item   не учитываются 
                пространственные характеристики видео;
        \item   не учитывается
                семантическая информация видео;
        \item   требуется настройка размера скользящих 
                окон при поиске аномалий;
        \item   требуется настройка коэффициента дискретизации 
                при построении индекса и поиску по индексу.
    \end{itemize}

\end{document}
